\newpage
\section{城轨车辆制造工艺概述}
在中车浦镇，我们重点了解了城轨车辆，包括地铁、市域车等的制造工艺，下面从转向架、车体、车辆总装这几个方面进行总结。

\subsection{转向架制造与组装工艺}
转向架是车辆最重要的部件之一，其结构之复杂、功能之重要使其成为制造工艺最复杂、精度要求最高的部件之一。\par{}

\subsubsection{轮对制造}
轮对是转向架的重要组成部分，由轴与轮固结成一个整体，因此轮对制造主要也分为轴和轮两部分。\par{}
对于轴，首先通过车床对其毛坯进行粗加工，再通过磨床对轴进行精加工与修正，最后做探伤检测。\par{}
对于轮，依旧是先对轮饼原料进行加工，再成形并去除其铁屑毛刺等，接着进入预组阶段，即在轮饼与轴接触的内圆上涂抹$\mathrm{MoS_2}$润滑脂方便组装。\par{}
预组之后进行正式压装。由于轴和轮的接触面呈现较为细微的锥度——靠近轴中点的部分宽、反之较窄——同时存在过盈配合，因此可以直接压紧轮与轴。压装好后还需进行测量，如测量轮位差等，防止误差过大影响舒适性与耐用性。\par{}
此外还需要进行装盘（轮盘式或轴盘式）、清洗、装降噪环等步骤，同时对轴和轮分别进行油漆，并在可能松动旋转的位置布置标记。\par{}
生产实习时，我们注意到地面上摆放着出口到印度班加罗尔的转向架的轮对。由于当地气候炎热潮湿，甲方要求在轮饼上和轴上都涂油漆；而国内一般只涂在轴上，较少给轮饼上漆。\par{}
轮对组装完毕后还要进行轴承压装、轴箱安装等，有单独的车间进行作业。由于该工序对温度较为敏感，因此为了防止热胀冷缩导致无法安装，必须把轮对与轴承存放至该密闭恒温车间并在同一温度下保温$8~\mathrm{h}$。\par{}

\subsubsection{构架焊接}
构架是转向架的主要承载、连接部件，其焊接质量决定了转向架的安全性，近年来浦镇公司采用人工智能、数字孪生等技术，将步骤繁琐的构架焊接产线进行了自动化改造，形成了构架焊接智能产线。\par{}
构架的焊接流程如下。首先由来料口进货并堆放至物料区，由自动化机器人进行点焊等装配工序；初步装配后由AGV运送至立库或焊接站，由7台焊接机器人进行焊接，此过程中只需要2人进行操作，焊接后进行磁粉探伤、尺寸检测、打磨等步骤，并进行整体焊接；期间会有AGV小车穿梭于机器与立库之间运送物料。如果设备发生异常，系统会第一时间通知一线工人，同时报警信息也会传送至中控室的数字孪生界面，以便进行修复。\par{}
图\ref{fig:goujia-welding}上半部分是构架焊接智能产线的大致布局，下半部分是该产线在数字孪生系统中的显示情况\sourcecite{2}。\par{}
\reportFigure{2}

\subsubsection{转向架整体组装}
转向架组装产线主要有4个工位，分别是节点加装与检测、挂件、落轮和静载/制动性能校验，呈U形排列，如图\ref{fig:bogie-u-line}：\par{}
\reportFigure{3}
第一步就是在构架上加装吊挂设备节点并进行检测。第二步指在转向架构架上加装对应部件，如制动缸的闸线，该步骤中构架有正位、反位两种放置方法，正位对应正常运行时转向架的上下方向，反位则反之，分别用于安装构架上下两边的部件。第三步将转向架安装并紧固到轮对上，浦镇厂中的城轨列车一般采用导柱式和转臂式定位。第四步则是做各种力学性能测试。\par{}
厂内大量采用智能设备，如智能扳手，可记录员工与对应工作所使用的扭力等相关信息；员工可通过类似手机的手持终端接收任务或上报情况；完成某些工序后还要进行AI检测，但只作最终确认，大部分检查还是按照传统人工“自检+互检+专检”的模式进行。\par{}
其中，AI辅助检测采用了转向架AI机器视觉外观检测技术，采用非接触式、智能化机器视觉技术，通过图片对比找差，识别标记、防锈油是否异常，是否有异物存在，螺栓有无遮挡等，构建多维度精准感知与决策的“AI大脑”。其搭载六自由度机械臂与地沟多向采集模组，实现检测死角全覆盖；采用3D视觉相机，同步触发采集多维度空间图像数据；植入深度学习算法，对海量图片进行网络训练生成预测模型，实时输出精确检测结论。该设备效益十分明显，直接实现了无人化检测（原需要两人），效率提升了$50\%$（$60~\min\to30~\min$），检出率提升了$14\%$（$85\%\to99\%$），并可以自动生成数字化报告，取代了手写纸质报告。\par{}

\subsection{车体制造工艺}
城轨车辆的车体制造工艺相比转向架而言较为简单，但仍然十分重要。对于列车而言，如果说转向架是行走的脚步、内装是血肉，那么车体就是轻量化骨架。\par{}

\subsubsection{车体分类}
不论是国铁动车组还是地铁、城轨车，头车和中间车的车体都有所不同：$\mathrm{T_c}$车有驾驶室、无牵引电机，而$\mathrm{M}$、$\mathrm{M_p}$车无驾驶室、有牵引电机。它们的底架等部件都有所不同，如底架按有无司机室划分为头部底架和中间车底架、车顶也分为头车用和中间车用。\par{}

\subsubsection{车体连接方式与工艺}
车体的连接方式、连接工艺也多种多样，如铆接、焊接等，焊接又分为电弧焊、摩擦焊、激光焊等。一般而言铆接的密封性不如焊接，因此车体大都采用焊接的方式，部分列车或列车的某些部分仍然采用铆接；底架使用铆钉连接至横梁，下挂电池、配重等；侧墙有的分块安装，有的拼成整块统一安装，各有优缺点；端墙上应预留贯通道的位置；司机室应配备有防爬器等装置，部分驾驶室的端部还需要配备逃生门等。\par{}
其中，最具有代表性的焊接工艺就是铝合金搅拌摩擦焊技术。该技术的原理是高速旋转的搅拌头与工件摩擦产生热量，使得材料热塑性化，搅拌头沿着待焊界面移动，促使热塑性材料在搅拌头作用下相互扩散、混合，进而实现固态连接。由于不直接依靠电弧加热，该技术具有绿色环保、高效率、变形小、强度高等优势。\par{}
同时，浦镇厂还开发了车体底架尺寸自动化检测技术、底架激光投影定位技术、铝合金氧化膜自动化激光清洗技术等，极大提升了生产效率与精度。\par{}

\subsubsection{车体制造流程}
一般而言，车辆的底架、车顶、侧墙、端墙被分别制造，成型后直接进行车体预组，再整车焊接、检查、交验。具体而言主要分为底架组装、侧墙组装、端角柱组装、端墙组装、车顶组装、车顶定位焊、预组尺寸检测、车体焊接与检验等步骤。由于浦镇厂的产品制造具有小批量、多品种的特点，焊接、检测等工序数字化、自动化的优势也就显现了出来。\par{}

\subsection{车辆总装}
车辆总装是在车体、转向架、涂装等前序工序完成后，将各系统、内装、外部部件和调试流程集成起来的最终制造环节，浦镇总装厂房附近布局如图\ref{fig:general-assembly-layout}所示。\par{}
\reportFigure{4}
总装分厂厂房分为$300~\mathrm{m}\times60~\mathrm{m}$或$300~\mathrm{m}\times42~\mathrm{m}$两种尺寸，总装工作主要在$42~\mathrm{m}$跨的厂房里的高站流水线上进行，共有14个工位，人员随作业任务在不同工位间调配，每个工位的工作内容不固定。参观当日，本人所在工位正准备对上海地铁21号线的第二批列车的某$\mathrm{M_p}$车进行$5~\mathrm{kV}$耐压测试。该车以黑色打底，白色、橙色、金色波浪线横贯车体，图案具有动感；整车采用内藏式车门，开门时门缩入墙体内；供电制式采用$\mathrm{DC}~1500~\mathrm{V}$受电弓受流，车辆外观与内装效果如图\ref{fig:l21-train}\sourcecite{3}所示。据工人师傅说，上海地铁要求实现的功能较为丰富，因此该型列车的设备较复杂。\par{}
\reportFigure{5}
同时，在$60~\mathrm{m}$跨厂房内，还存放着杭州地铁列车、穗莞深城际列车、出口至沙特阿拉伯的跨座式单轨列车和出口至格鲁吉亚第比利斯的地铁列车等，可见浦镇厂产品的多样性。

